\documentclass[french,a4paper]{article}

\usepackage[utf8]{inputenc}
\usepackage[T1]{fontenc}
\usepackage{lmodern}
\usepackage{geometry}
\usepackage{graphicx}
\usepackage{babel}
\usepackage{amsmath}
\usepackage{amsthm}
\usepackage{amssymb}
\usepackage{hyperref}
\usepackage{stmaryrd}
\usepackage{enumitem}
\usepackage{tcolorbox}
\usepackage{dsfont}
\usepackage{multirow}
\usepackage{etoc}
\usepackage{titlesec}
\usepackage{esint}
\usepackage{cancel}
\usepackage{mathdots}

\setlist[itemize]{label=$\bullet$}


\renewcommand{\P}{\mathbb{P}}
\newcommand{\N}{\mathbb{N}}
\newcommand{\Z}{\mathbb{Z}}
\newcommand{\Q}{\mathbb{Q}}
\newcommand{\R}{\mathbb{R}}
\newcommand{\C}{\mathbb{C}}
\newcommand{\K}{\mathbb{K}}
\renewcommand{\L}{\mathbb{L}}
\newcommand{\U}{\mathbb{U}}
\newcommand{\st}{^*}
\newcommand{\tq}{\middle|}
\newcommand{\inv}{^{-1}}
\newcommand{\E}{\mathbb{E}}
\newcommand{\F}{\mathbb{F}}
\newcommand{\V}{\mathbb{V}}
\renewcommand{\ker}{\text{Ker}}
\newcommand{\im}{\text{Im}}
\newcommand{\card}{\text{Card}}
\newcommand{\Tr}{\text{Tr}}

\theoremstyle{plain}
\newtheorem{thm}{Théorème}
\newtheorem*{ind}{Indication}

\theoremstyle{definition}
\newtheorem{dfn}{Definition}

\theoremstyle{remark}
\newtheorem*{rmq}{Remarque}
\newtheorem*{app}{Application}

\newtcolorbox[auto counter]{ex}[2][]{colback=white, colframe=green!50!white, coltitle=black, fonttitle=\bfseries, title={Exercice~\thetcbcounter~: #2}, after title={\hfill #1}}

\title{TD Intégrales à paramètres}

\begin{document}
\maketitle

Dans plusieurs exercices, on aura besoin du résultat suivant : $$\int_{0}^{+\infty}e^{-t^2}\ dt=\frac{\sqrt{\pi}}{2}$$

\section{Suites et séries d'intégrales}
\begin{ex}[Moulin.IP.01]{$\circ$}
	Soit $f$ intégrable sur $\R$. Montrer la convergence de l'intégrale $$I_n=\int_{-\infty}^{+\infty}f(x)e^{-x^2/n^2}\ dx$$ et déterminer sa limite lorsque $n$ tend vers $+\infty$.
\end{ex}
\begin{proof}
	Pour l'existence, on a de la continuité par morceaux et un $O(f)$ en $-\infty$ et en $+\infty$. La limite simple est la fonction nulle sauf en $0$ où c'est $f(0)$. On a une domination par $\lvert f\rvert$ qui est intégrable donc la limite est nulle.
\end{proof}

\begin{ex}[Moulin.IP.02]{}
	Montrer que $$I_n=\int_{1}^{+\infty}e^{-x^n}\ dx$$ tend vers $0$ et en trouver un équivalent en utilisant le changement de variable $t=x^n$.
\end{ex}
\begin{proof}
	La limite simple est la fonction nulle partout sauf en $1$ où elle vaut $e^{-1}$. On a une domination par $e^{-x}$ qui intégrable donc la limite est nulle d'après le théorème de convergence dominée. Avec le changement de variable on trouve : $$I_n=\frac{1}{n}\int_{1}^{+\infty}\frac{t^{1/n}e^{-t}}{t}\ dt$$ On a cette fois encore une domination par $e^{-t}$ qui nous fournit avec le théorème de convergence dominée : $$I_n\sim\frac{1}{n}\int_{1}^{+\infty}\frac{e^{-t}}{t}\ dt$$ (cette intégrale est bien non nulle car la fonction est continue positive non identiquement nulle).
\end{proof}

\begin{ex}[Moulin.IP.03]{Recherche de développement asymptotique}
	On définit : $$I_n=\int_{0}^{1}\frac{x^n}{1+x^n}\ dt$$
	\begin{enumerate}
		\item Déterminer un équivalent de $(I_n)$.
		\item Donner un développement asymptotique à deux termes de $(I_n)$.
	\end{enumerate}
\end{ex}
\begin{proof}
	On fait comme dans l'exercice précédent.
	\begin{enumerate}
		\item Le changement de variable $u=x^n$ donne : $$I_n=\frac{1}{n}\int_{0}^{1}\frac{u^{1/n}}{1+u}\ du$$ Le théorème de convergence dominée, avec une domination par $\displaystyle\frac{1}{1+u}$ donne alors : $$I_n\sim\frac{1}{n}\int_{0}^{1}\frac{du}{1+u}=\frac{\ln(2)}{n}$$
		\item Puis : $$I_n-\frac{\ln(2)}{n}=\frac{1}{n}\int_{0}^{1}\frac{u^{1/n}-1}{1+u}\ du=\frac{1}{n^2}\int_{0}^{1}\frac{n(u^{1/n}-1)}{1+u}\ du$$ On a pour cette nouvelle intégrale une convergence simple vers $\displaystyle\frac{\ln(u)}{1+u}$ et une domination par $\displaystyle\frac{-\ln(u)}{1+u}$ qui est intégrable. Par convergence dominée, on a : $$\int_{0}^{1}\frac{n(u^{1/n}-1)}{1+u}\ du\xrightarrow[n\to+\infty]{}\int_{0}^{1}\frac{\ln(u)}{1+u}\ du=\sum_{n=1}^{+\infty}\frac{(-1)^n}{n^2}=-\frac{\pi^2}{12}$$ (l'avant-dernière égalité se prouver par intégration terme à terme, puis on prouve la dernière en sortant les termes pairs). On a donc le développement asymptotique : $$I_n=\frac{\ln(2)}{n}-\frac{\pi^2}{12n^2}+o\left(\frac{1}{n^2}\right)$$
	\end{enumerate}
\end{proof}

\begin{ex}[Moulin.IP.04]{}
	On définit $$I_n=\int_{-\infty}^{+\infty}\left(1+\frac{x^2}{n}\right)^{-n}\ dx$$ Montrer que $(I_n)_{n\in\N}$ converge et admet pour limite $$\int_{-\infty}^{+\infty}e^{-x^2}\ dx$$
\end{ex}
\begin{proof}
	On vérifie que la limite simple est $e^{-x^2}$ sans trop de problème. Pour la domination, cela se corse. Il faut étudier $g:t\mapsto\displaystyle e^{-t\ln\left(1+\frac{x^2}{t}\right)}$ pour $x$ fixé lorsque $t\in\R_+\st$. En la dérivant et en utilisant l'inégalité : $$1-\frac{1}{t}\leqslant\ln(t)\leqslant t-1$$ on prouver qu'elle est décroissante. On en déduit la décroissance de $(f_n(x))$ et donc une domination par $f_1(x)=\displaystyle\frac{1}{1+x^2}$ qui est intégrable. Le théorème de convergence dominée conclut.
\end{proof}

\begin{ex}[Moulin.IP.05]{Théorème de convergence encadrée}
	Soit $(f_n)_{n\in\N}$ une suite de fonctions continues par morceaux sur un intervalle $I$ et convergeant simplement vers $f$ continue par morceaux sur $I$. On suppose qu'il existe deux fonctions $\varphi$ et $\psi$ continues par morceaux sur $I$ telles que pour tout $n\in\N$, $\varphi\leqslant f_n\leqslant\psi$, et dont les intégrales sur $I$ convergent. Montrer que les intégrales des $f_n$ et de $f$ convergent sur $I$, et que l'on a : $$\int_I f=\lim_{n\to+\infty} \int_I f_n$$
\end{ex}

\begin{ex}[Moulin.IP.06]{$\star$}
	Soit $(f_n)_{n\in\N}$ une suite de fonctions positives intégrables sur un intervalle $I$ convergeant simplement vers uen fonction $f$ continue par morceaux non intégrable. Montrer : $$\lim_{n\to+\infty}\int_I f_n=+\infty$$ On pourra considérer $\inf(f_n,f)$.
\end{ex}

\begin{ex}[Moulin.IP.07]{}
	Convergence et limite lorsque $n$ tend vers $+\infty$ de $I_n=\displaystyle\int_{0}^{1}x^{-1/n}\left(1+\frac{x}{n}\right)^{-n}\ dx$
\end{ex}

\begin{ex}[Moulin.IP.08]{}
	Convergence et limite lorsque $n$ tend vers $+\infty$ de $J_n=\displaystyle\int_{0}^{n}x^{-1/n}\left(1-\frac{x}{n}\right)^{n}\ dx$
\end{ex}

\begin{ex}[Moulin.IP.09]{}
	Montrer l'égalité $$\int_{0}^{+\infty}\frac{\sin(x)}{e^x-1}\ dx=\sum_{n=1}^{+\infty}\frac{1}{1+n^2}$$
\end{ex}
\begin{proof}
	Écrire $$\frac{\sin(x)}{e^x-1}=\sin(x)\frac{e^{-x}}{1-e^{-x}}=\sin(x)\sum_{n=1}^{+\infty e^{-nx}}$$ On pose $u_n(x)=\sin(x)e^{-nx}$. On montre que $\displaystyle\int_{0}^{+\infty}\lvert u_n(x)\rvert\ dx$ est sommable en découpant entre $0$ et $1$ pour utiliser la majoration classique de $\lvert\sin(x)\rvert$ par $\lvert x\rvert$ puis entre $1$ et $+\infty$ où on peut simplement calculer l'intégrale. On majore alors en valeur absolue par $$\frac{1}{n^2}(1-e^{-n})\sim\frac{1}{n^2}$$ qui est sommable (des termes se sont simplifiés). On peut donc intégrer terme à terme. Ensuite, on calcule les intégrales des $u_n$ en passant par une partie imaginaire d'une exponentielle complexe pour obtenir le résultat souhaité.
\end{proof}
\begin{ex}[Moulin.IP.10]{Géométrique !}
	Soit $p$ et $q$ deux réels strictement positifs. Montrer l'égalité $$\int_{0}^{1}\frac{x^{p-1}}{1+x^q}=\sum_{n=0}^{+\infty}\frac{(-1)^n}{nq+p}$$
\end{ex}
\begin{proof}
	Si $x\in[0,1]$, on écrit : $$\frac{x^{p-1}}{1+x^q}=\sum_{n=0}^{+\infty}\underbrace{(-1)^nx^{nq+p-1}}_{:=u_n(x)}$$ Ensuite, on veut intégrer terme à terme, mais $$\int_{0}^{+\infty}u_n(x)\ dx=\frac{(-1)^n}{nq+p} \ \ \text{et}\ \ \int_{0}^{+\infty}\lvert u_n(x)\rvert\ dx=\frac{1}{nq+p}$$ donc le théorème d'intégration terme à terme ne s'applique pas. Cependant, grâce au théorème des séries alternées, on a, puisque $\lvert u_n(x)\rvert$ décroît : $$\left\lvert\int_{0}^{1}R_{n-1}(x)\ dx\right\rvert\leqslant\int_{0}^{1}\lvert R_{n-1}(x)\rvert\ dx\leqslant\int_{0}^{1}\lvert u_n(x)\rvert dx=\frac{1}{nq+p}\xrightarrow[n\to+\infty]{}0$$ donc $$\int_{0}^{1}R_{n-1}(x)\ dx\xrightarrow[n\to+\infty]{}0$$ En conséquence, on peut tout de même intervertir somme et intégrale, ce qui fournit le résultat.
\end{proof}
\begin{ex}[Moulin.IP.11]{$\Delta$ Sophomore's dream}
	Montrer que la fonction $x\mapsto x^x$ est prolongeable en une fonction $f$ continue sur $\R_+$. Montrer que $$\int_{0}^{1}f(x)\ dx=\sum_{n=1}^{+\infty}\frac{(-1)^{n-1}}{n^n}$$
\end{ex}
\begin{proof}
	Avec un DL, on voit que le seul prolongement convenable est le prolongement par $1$ en $0$. Ensuite, développer $f$ en série avec le développement en série de l'exponentielle puis calculer : $$\int_{0}^{1}\frac{(x\ln(x))^n}{n!}\ dx=\frac{(-1)^n}{(n+1)^{n+1}}$$ Pour cela, il est préférable de définir $I_{n,p}$ avec une puissance sur le $\ln$ et l'autre paramètre pour la factorielle puis de faire une récurrence sur le nombre de $\ln$ en intégrant par parties. Or, puisque le signe de ces termes est constant, on a $$\left\lvert \int_{0}^{1}\frac{(x\ln(x))^n}{n!}\ dx\right\rvert=\int_{0}^{1}\left\lvert\frac{(x\ln(x))^n}{n!}\right\rvert\ dx=\frac{1}{(n+1)^{n+1}}$$ qui est sommable car $O\displaystyle\left(\frac{1}{n^2}\right)$. Ensuite, intervertir somme et intégrale pour obtenir le résultat.
\end{proof}

\begin{ex}[Moulin.IP.13]{Beurk}
	Montrer que $$f:x\mapsto\sum_{n=1}^{+\infty}\frac{(-1)^{n+1}\ln(n+1)}{n(n+1)^{x+1}}$$ est intégrable sur $\R_+$ et calculer son intégrale.
\end{ex}
\begin{proof}
	On montre facilement que $u_n(x)$ est intégrable. Ensuite, on considère : $$f:t\mapsto\frac{\ln(t+1)}{t(t+1)^{x+1}}$$ à $x$ fixé que l'on dérive. En utilisant l'inégalité $\ln(t)\geqslant1-\displaystyle\frac{1}{t}$, on prouve, en la dérivant, que $f$ est décroissante et donc que $(\lvert u_n(x)\rvert)$ est décroissante. On calcule : $$\int_{0}^{+\infty}\lvert u_n(x)\rvert\ dx=\frac{1}{n(n+1)}$$ et on peut alors intégrer terme à terme. On intervertit somme et intégrale et, en écrivant une somme télescopique (après avoir sorti un terme), on prouve que l'intégrale vaut $2\ln(2)-1$.
\end{proof}

\begin{ex}[Moulin.IP.14]{}
	\begin{enumerate}
		\item Pour $x>0$, on définit : $$H(x)=\int_{0}^{+\infty}\lvert \sin(y)\rvert e^{-xy}\ dy$$ Montrer que cela a bien un sens, calculer $H(x)$ puis déterminer un équivalent en $+\infty$ de $H$.
		\item Calculer, pour $z>0$ et $x$ réel, l'intégrale $\displaystyle\int_{0}^{+\infty}\sin(xy)e^{-zy}\ dy$
		\item Montrer, pour tout $t\in\R$ et $x>-1$, l'égalité : $$\int_{0}^{+\infty}e^{-xy}\frac{\sin(ty)}{e^y-1}\ dy=\sum_{n=1}^{+\infty}\frac{t}{t^2+(n+x)^2}$$
	\end{enumerate}
\end{ex}

\begin{ex}[Moulin.IP.15]{$\star$}
	Soit $a$ un nombre réel. Déterminer un équivalent lorsque $n$ tend vers $+\infty$ de $$I_n(a)=\int_{-\infty}^{a}\frac{x^2e^{-nx^2}}{1+x^2}\ dx$$ On discutera selon le signe de $a$.
\end{ex}

\begin{ex}[Moulin.IP.16]{$\star\star$}
	Déterminer un équivalent de $$u_n=\displaystyle\sum_{k=n}^{+\infty}\frac{n^k}{k!}$$ On pourra utiliser la formule de Stirling pour se ramener à étudier $\displaystyle\int_{n}^{+\infty}\frac{(en)^x}{x^x\sqrt{2\pi x}}\ dx$ puis effectuer le changement de variable $x=n+u$ dans cette intégrale.
\end{ex}

\begin{ex}[Moulin.IP.17]{Théorème de convergence monotone}
	Soit $(f_n)$ une suite croissante de fonctions intégrables convergeant simplement vers une fonction $f$ continue par morceaux. Montrer que $f$ est intégrable si, et seulement si, la suite $\left(\displaystyle\int_{I}f_n\right)$ est majorée et qu'alors on a : $$\int_{I}f=\lim_{n\to+\infty}\int_I f_n$$
\end{ex}
\begin{proof}
	Historiquement, pour démontrer le théorème de convergence dominée, on a commencé par démontrer celui-ci. On considère $(f_{n+1}-f_n)$. Alors $\sum(f_{n+1}-f_n)$ CVS vers $f-f_0$ CPM. Alors, par croissance, tous les termes qui suivent existent dans $\overline{\R_+}$ : $$\int_I (f-f_0)=\sum_{n=0}^{+\infty}\int_I(f_{n+1}-f_n)=\sum_{n=0}^{+\infty}\left(\int_I f_{n+1}-\int_I f_n\right)=\lim_{n\to+\infty}\left(\int_I f_n\right)-\int_If_0$$ Et alors on raisonne par équivalences en remarquant que $f-f_0$ est positive donc la convergence de son intégrale équivaut à son intégrabilité.
\end{proof}
\section{Intégrales dépendant d'un paramètre réel}
\begin{ex}[Moulin.IP.18]{$\circ$}
	Définition, continuité, dérivabilité et calcul de $$f(x)=\int_{0}^{+\infty}\frac{\sin(xt)}{t}e^{-t}\ dt$$
\end{ex}
\begin{proof}
	$O(e^{-t})$ pour la définition. pour la continuité, une domination par $ae^{-t}$ sur tout $[-a,a]$ pour $a>0$ donne la continuité sur $\R$. Pour la dérivabilité, on domine $$\left\lvert\frac{\partial g}{\partial x}(x,t)\right\rvert=\lvert \cos(xt)\rvert e^{-t}\leqslant e^{-t}$$ En calculant, on trouve ensuite que $$f'(x)=\frac{1}{1+x^2}$$ (passer par des exponentielles complexes). Or, on voit immédiatement que $f(0)=0$ et alors $f(x)=\arctan(x)$.
\end{proof}
\begin{rmq}
	On peut en déduire la valeur de $$\int_{0}^{+\infty}\frac{\sin(u)}{u}\ du$$ En effet, pour $x>0$, on a $$f(x)=\int_{0}^{+\infty}\frac{\sin(u)}{u}e^{-u/x}\ du$$ Ensuite, on pose $$f_n(x)=\int_{0}^{n\pi}\frac{\sin(u)}{u}e^{-u/x}\ du$$ On vaut alors utiliser le théorème de la double limite. Pour la convergence uniforme par rapport à $x$ lorsque $n\to+\infty$, on utilise des séries alternées en découpant les intégrales. Et le théorème de convergence dominée donne $$f_n(x)\xrightarrow[x\to+\infty]{}\int_{0}^{n\pi}\frac{\sin(u)}{u}\ du$$ On en déduit avec l'expression de $f$ que $$\int_{0}^{+\infty}\frac{\sin(u)}{u}\ du=\frac{\pi}{2}$$
\end{rmq}
\begin{ex}[Moulin.IP.19]{Une étude}
	Domaine de définition et calcul de : $$f:x\mapsto\int_{0}^{+\infty}\frac{1-\cos(t)}{t^2}e^{-xt}\ dt$$
\end{ex}
\begin{proof}
	On dérive deux fois (en dominant sur $[a,+\infty[$) puis on calcule pour obtenir : $$f''(x)=\frac{1}{x}-\frac{x}{x^2+1}$$ Ensuite, quand on primitive une première fois, la constante est nulle avec la limite en $+\infty$ puis en primitivant une deuxième fois, la limite en $+\infty$ donne $\pi/2$ pour la deuxième constante. On obtient : $$f(x)=x\ln(x)-\frac{1}{2}\ln(1+x^2)-\arctan(x)+\frac{\pi}{2}$$
\end{proof}
\begin{rmq}
	Par continuité, on en déduit que $$\int_{0}^{+\infty}\frac{1-\cos(t)}{t^2}\ dt=\frac{\pi}{2}$$
\end{rmq}
\begin{ex}[Moulin.IP.20]{$\Delta$ Prolongements $\mathcal{C}^\infty$}
	\begin{enumerate}
		\item Montrer que si $f$ est de classe $\mathcal{C}^\infty$ sur $\R$ et $f(0)=\ldots=f^{(n-1)}(0)=0$, alors il existe une fonction de classe $\mathcal{C}^\infty$ telle que $$\forall x\in\R,\ f(x)=x^ng(x)$$ On pourra utiliser une formule de Taylor.
		\item En déduire que si $f$ est de classe $\mathcal{C}^\infty$ sur $\R$ et si $P$ est un polynôme tel que la fonction $t\mapsto f(t)/P(t)$ admette un prolongement continu, alors ce prolongement est de classe $\mathcal{C}^\infty$ sur $\R$.
	\end{enumerate}
\end{ex}
\begin{proof}
	\begin{enumerate}
		\item Suivre l'indication et utiliser la formule de Taylor avec reste intégral. Montrer que l'intégrale, lorsqu'on a sorti $x^n$ est de classe $\mathcal{C}^\infty$ en dominant les dérivées sur tout segment.
		\item Se ramener en $0$ par translation et appliquer la question précédente au numérateur.
	\end{enumerate}
\end{proof}

\begin{ex}[Moulin.IP.21]{}
	Soit $I$ un intervalle de $\R$. Pour $f\in\mathcal{C}^0(I,\C)$ et $(a,b)\in I^2$ tel que $a\ne b$, on pose : $$I_f(a,b)=\int_{a}^{b}\frac{f(x)}{\sqrt{(x-a)(b-x)}}\ dx$$
	\begin{enumerate}
		\item Montrer la convergence de l'intégrale $I_f(a,b)$ et la calculer lorsque $f$ est une fonction constante.
		\item Montrer que $I_f$ est continue sur $I^2\setminus\{(x,x)\mid x\in I\}$. Peut-on la prolonger par continuité sur $I^2$ ?
	\end{enumerate}
\end{ex}

\begin{ex}[Moulin.IP.22]{}
	Étudier la fonction définie par : $$F(x)=\int_{0}^{+\infty}\frac{dt}{\sqrt{1+t^2}\sqrt{x^2+t^2}}$$ (ensemble de définition, parité, relation entre $F(x)$ et $F(1/x)$, régularité, limites aux bornes, $\star$ équivalents aux bornes)
\end{ex}

\begin{ex}[Moulin.IP.23]{}
	Continuité et dérivabilité de $$F:x\mapsto\int_{0}^{+\infty}\frac{e^{-xt}}{1+t^2}\ dt$$ Est-elle dérivable en $0$ ?
\end{ex}
\begin{proof}
	$F$ est définie sur $\R_+$. $F$ est continue sur $\R_+$ en dominant par : $$t\mapsto\frac{1}{1+t^2}$$ $F$ est continue sur $\R_+\st$ en dominant $t\mapsto\displaystyle\frac{-te^{-xt}}{1+t^2}$ par $$t\mapsto\frac{te^{-at}}{1+t^2}$$ pour $x\in[a,+\infty[$. Le théorème de dérivabilité ne permet pas de conclure sur la dérivabilité ou non en $0$. On calcule : $$\frac{F(x)-F(0)}{x}=\int_{0}^{+\infty}\frac{e^{-xt}-1}{x(1+t^2)}dt=\int_{0}^{+\infty}\frac{e^{-u}-1}{x^2\left(1+\frac{u^2}{x^2}\right)}du=\int_{0}^{+\infty}\frac{e^{-u}-1}{x^2+u^2}du$$ Il semblerait alors que $F$ ne soit pas dérivable en $0$ car cette intégrale semble diverger vers $-\infty$ quand $x$ tend vers $0$ (à vérifier).
\end{proof}

\begin{ex}[Moulin.IP.24]{}
	Soit $f$ une fonction continue sur $[0,1]$ et à valeurs strictement positives. En utilisant $$G(x)=\int_{0}^{1}\exp(x\ln(f(t)))\ dt$$ trouver la limite lorsque $x$ tend vers $0$ de $$F(x)=\left(\int_{0}^{1}f(t)^x\right)^{1/x}$$
\end{ex}

\begin{ex}[Moulin.IP.25]{}
	Pour tout $x\in\R$, on pose : $f(x)=\displaystyle\int_{0}^{+\infty}\frac{\sin(xt)}{(1+t^2)\sqrt{t}}\ dt$
	\begin{enumerate}
		\item Montrer que $f$ est de classe $\mathcal{C}^1$ sur $\R$. Calculer $f'(0)$.
		\item Montrer que $f$ est de classe $\mathcal{C}^\infty$ sur $\R\st$. Est-elle deux fois dérivable en $0$ ?
		\item Montrer que $f$ est positive sur $\R_+$.
	\end{enumerate}
\end{ex}

\begin{ex}[Moulin.IP.26]{}
	Montrer que dans le théorème de dérivation sous le signe intégral de $\displaystyle\int_If(x,t)\ dt$, on peut remplacer l'hypothèse d'intégrabilité, pour tout $x$, de $t\mapsto f(x,t)$ par la convergence d'une intégrale $\displaystyle\int_I f(a,t)\ dt$ pour un certain $a$.
\end{ex}

\begin{ex}[Châteaux]{$\star$ Théorème de Bohr-Mollerup}
	On rappelle la définition de $\Gamma:x\mapsto\displaystyle\int_{0}^{+\infty}t^{x-1}e^{-t}\ dt$ pour $x>0$. On veut démontrer le théorème de Bohr-Mollerup, qui stipule que $\Gamma$ est l'unique solution de l'équation fonctionnelle : $$f(x+1)=xf(x)$$ avec $f:\R_+\st\to\R_+\st$ qui vérifie de plus $f(1)=1$ et $\ln\circ f$ convexe.
	\begin{enumerate}
		\item Vérifier que $\Gamma$ est solution.
		\item Démontrer que pour tout $x>0$ : $\displaystyle\frac{n^xn!}{x(x+1)\cdots(x+n)}\xrightarrow[n\to+\infty]{}\Gamma(x)$
		\item Soit $f$ vérifiant l'équation fonctionnelle.
		\begin{enumerate}[label=\alph*.]
			\item Montrer : $\displaystyle\forall x\in]0,1[,\ \ln(n)\leqslant\frac{1}{x}\ln\left(\frac{f(n+1+x)}{f(n+1)}\right)\leqslant\ln(n+1)$
			\item En déduire : $$\forall x\in]0,1[,\ \displaystyle\frac{n^xn!}{x(x+1)\cdots(x+n)}\leqslant f(x)\leqslant\left(1+\frac{1}{n}\right)^x\frac{n^xn!}{x(x+1)\cdots(x+n)}$$
			\item Conclure que $f=\Gamma$.
		\end{enumerate}
	\end{enumerate}
\end{ex}
\begin{proof}
	\begin{enumerate}
		\item IPP pour la relation fonctionnelle. Calcul immédiat pour $f(1)$. Pour montrer que $\ln\circ\Gamma$ est convexe, dériver deux fois $\Gamma$ et utiliser l'inégalité de Cauchy-Schwarz.
		\item Utiliser le théorème de convergence dominée avec $$I_n=\int_{0}^{n}t^{x-1}\left(1-\frac{t}{n}\right)^{n}\ dt=\frac{n!n^{x+1}}{x(x+1)\cdots(x+n)n^x}$$ qui, complétée par $0$, tend vers $\Gamma(x)$ (la domination par $t^{x-1}e^{-t}$ est facile). 
		\item \begin{enumerate}[label=\alph*.]
			\item Utiliser le lemme des trois pentes puis le fait que $f(n+1)=n!$ (par récurrence avec l'hypothèse fonctionnelle).
			\item Partir de la sous-question précédente, multiplier par $x$, passer à l'exponentielle, multiplier par $f(n+1)=n!$ puis utiliser la propriété fonctionnelle pour montrer que $f(n+1+x)=x(x+1)\cdots(x+n)f(x)$.
			\item On en déduit par passage à la limite $\forall x\in]0,1[,\ f(x)=\Gamma(x)$ et on a $f(1)=\Gamma(1)=1$ par hypothèse. Comme $f$ et $\Gamma$ sont entièrement déterminée par leurs valeurs sur $]0,1[$ au vu de la propriété fonctionnelle $f(x+1)=xf(x)$, on en déduit que $f=\Gamma$.
		\end{enumerate}
	\end{enumerate}
\end{proof}

\begin{ex}[Châteaux]{Intégrale de Fresnel}
	\begin{enumerate}
		\item Montrer que $$\int_\R\frac{x^2}{1+x^4}dx=\int_\R\frac{1}{1+x^4}dx$$ et calculer cette intégrale.
		\item On pose, pour tout $t\in\R_+$, $F(t)=\displaystyle\int_{0}^{+\infty}\frac{\exp(-(x^2+i)t^2)}{x^2+i}\ dx$. Montrer que $F$ est continue et étudier sa limite en $+\infty$.
		\item Montrer que $F$ est de classe $\mathcal{C}^1$ sur $\R_+\st$ et calculer sa dérivée.
		\item Montrer que $\displaystyle\int_{0}^{+\infty}e^{it^2}\ dt$ converge et calculer sa valeur.
	\end{enumerate}
\end{ex}

\begin{ex}[Moulin.IP.27]{Intégrales de Fresnel}
	Pour $(x,y)\in\R_+\st\times\R$, on pose $$G(x,y)=\int_{0}^{+\infty}\frac{e^{-xt}e^{iyt}}{\sqrt{t}}\ dt$$
	\begin{enumerate}
		\item Montrer que $G$ est continue sur $\R_+\st\times\R$.
		\item Calculer $G(1,0)$.
		\item Exprimer $G(x,y)$ à l'aide de $G(1,z)$ pour un $z$ bien choisi.
		\item Montrer que $H:z\mapsto G(1,z)$ est dérivable sur $\R$ et que $$\forall z\in\R,\ H'(z)=\frac{i}{2(1-iz)}H(z)$$
		\item $\star$ En déduire la valeur de $H(z)$ pour tout $z\in\R$ puis celle de $G(x,y)$ pour tout $(x,y)\in\R_+\st\times\R$.
		\item Montrer que $G(0,y)$ a un sens pour tout $y>0$.
		\item $\star$ Montrer que $\displaystyle\lim_{x\to 0^+}G(x,1)=G(0,1)$. En déduire la convergence et la valeur de $$I=\int_{0}^{+\infty}\cos(t^2)\ dt \ \ \text{et} \ \ J=\int_{0}^{+\infty}\sin(t^2)\ dt$$
	\end{enumerate}
\end{ex}
\begin{proof}
	Les questions 5 et 7 sont difficiles.
	\begin{enumerate}
		\item  Dominer par $\displaystyle\frac{e^{-\varepsilon t}}{\sqrt{t}}$ sur $[\varepsilon,+\infty[\times\R$ et utiliser le théorème de continuité.
		\item Calculer en utilisant l'intégrale de Gauss ce qui donne avec un changement de variable $$G(1,0)=\sqrt{\pi}$$
		\item Le changement de variable $u=xt$ fournit : $$G(x,y)=\frac{1}{\sqrt{x}}G\left(1,\frac{y}{x}\right)$$
		\item Dériver avec le théorème. La domination ne pose pas de problème. Faire ensuite une IPP sur l'intégrale dérivée pour obtenir le résultat.
		\item Résoudre \textbf{proprement} (pas de $\ln$ complexe !) cette équation différentielle pour obtenir : $$\forall z\in\R,\ H(z)=\sqrt{\pi}\frac{1}{\left(1+z^2\right)^{1/4}}e^{\frac{i}{2}\arctan(z)}$$
		Avec le résultat de la question 3, on en déduit que : $$G(x,y)=\sqrt{\frac{\pi}{x}}{\left(1+\left(\frac{y}{x}\right)^2\right)^{-1/4}}e^{\frac{i}{2}\arctan\left(\frac{y}{x}\right)}$$
		\item On l'intégrabilité en $0$. La convergence en $+\infty$ s'établit par une IPP.
		\item Écrire
		\begin{align*}
			\left\lvert G(x,1)-\int_{1/n}^{n}\frac{e^{-xt}e^{it}}{\sqrt{t}}\ dt\right\rvert&\leqslant\int_{0}^{1/n}\frac{dt}{\sqrt{t}}+\left\lvert\left[\frac{e^{(i-x)t}}{(i-x)\sqrt{t}}\right]_n^{+\infty} -\int_{n}^{+\infty}\frac{-e^{(i-x)t}}{2(i-x)t\sqrt{t}}\right\rvert \\
			&\leqslant \frac{2}{\sqrt{n}}+\frac{1}{\sqrt{n}}\left\lvert\frac{x+i}{x^2+1}\right\rvert+\left\lvert\frac{x+i}{x^2+1}\right\rvert\int_{n}^{+\infty}\frac{dt}{2t\sqrt{t}}\\
			&\leqslant\frac{2}{\sqrt{n}}+\frac32\frac{1}{\sqrt{n}}+\frac32\frac{1}{\sqrt{n}}\\
			&=\underbrace{\frac{5}{\sqrt{n}}}_{\text{ind. de }x}\xrightarrow[n\to+\infty]{}0
		\end{align*} donc on a une convergence uniforme sur $\R_+$. Or, les $$x\mapsto \int_{1/n}^{n}\frac{e^{-xt}e^{it}}{\sqrt{t}}\ dt$$ sont continues sur $\R_+$ (dominer) donc $x\mapsto G(x,1)$ aussi donc on a le résultat souhaité. Ensuite, on utilise l'expression de la question 5 et la continuité pour écrire : $$\sqrt{\pi}e^{i\pi/4}=G(0,1)=2\int_{0}^{+\infty}e^{iu^2}\ du$$ après le changement de variable $u=\sqrt{t}$. En passant aux parties réelle et imaginaire, on obtient la convergence de $I$ et $J$ ainsi que $$I=J=\frac{1}{2}\sqrt{\frac{\pi}{2}}$$
	\end{enumerate}
\end{proof}

\begin{ex}[Moulin.IP.28]{$\star$ Fonctions d'Airy}
	Montrer que $F:x\mapsto\displaystyle\int_{-\infty}^{+\infty}\exp\left(i(tx+t^3/3)\right)\ dt$ est de classe $\mathcal{C}^1$ sur $\R$.
\end{ex}

\begin{ex}[Moulin.IP.29]{}
	Pour tout $x>0$, on pose : $$f(x)=\int_{0}^{+\infty}\frac{t-\lfloor t\rfloor}{t(t+x)}\ dt$$
	\begin{enumerate}
		\item Montrer que $f$ est de classe $\mathcal{C}^1$ sur $\R_+\st$ et en déterminer un équivalent au voisinage de $0^+$.
		\item $\star\star$ Déterminer un équivalent de $f$ au voisinage de $+\infty$.
	\end{enumerate}
\end{ex}

\begin{ex}[Moulin.IP.30]{}
	Soit $f$ une fonction réelle continue bornée et intégrable sur $\R_+$.
	\begin{enumerate}
		\item Quels sont les nombres complexes $z$ non nuls pour lesquels l'intégrale $$F(z)=\frac{1}{\pi}\int_{0}^{+\infty}\frac{f(t)}{t+z}\ dt$$ converge ?
		\item Montrer que pour tout $x\in\R_+\st$, on a : $$f(x)=\lim_{\varepsilon\to0^+}\im F(-x-i\varepsilon)$$ Et si $f$ est seulement continue par morceaux ?
	\end{enumerate}
\end{ex}
\begin{proof}
	\begin{enumerate}
		\item Si $z\in\R_-\st$, cela n'est pas défini. Pour $z=0$, on ne connaît pas \textit{a priori} le comportement de $\displaystyle\frac{f(t)}{t}$ au voisinage de $0$. Et si $z\in\C\setminus\R_-$, on obtient la bonne définition avec la continuité et un $O(f)$.
		\item Calculer $F(-x-i\varepsilon)$ en multipliant l'intégrande par le complexe conjugué du dénominateur. On obtient : $$\im F(-x-i\varepsilon)=\frac{1}{\pi}\int_{-\infty}^{+\infty}h(\varepsilon, u)\ du$$ avec $h(\varepsilon,u) =\frac{f(x+\varepsilon u)}{u^2+1}$ si $u\in[-\frac{x}{\varepsilon},+\infty[$ et $0$ sinon. On domine par $\displaystyle\frac{M}{u^2+1}$ avec $f$ bornée par $M$ et le théorème de passage à la limite sous l'intégrale fournit $$\im F(-x-i\varepsilon)\xrightarrow[\varepsilon\to 0^+]{}\frac{1}{\pi}f(x)\left(\frac{\pi}{2}-\frac{-\pi}{2}\right)=f(x)$$ Si $f$ est CPM, on peut prouver de même que : $$\im F(-x-i\varepsilon)\xrightarrow[\varepsilon\to 0^+]{}\frac12\left(f(x^-)+f(x^+)\right)$$
	\end{enumerate}
\end{proof}
\begin{ex}[Châteaux]{Cas où le théorème de dérivation ne fonctionne pas} Montrer que $$F:x\mapsto\int_{0}^{+\infty}\frac{\sin(xt)}{x(1+t^2)} \ dt$$ est de classe $\mathcal{C}^2$.
\end{ex} 
\begin{proof}
	On commence par utiliser le théorème de dérivation pour montrer qu'elle est $\mathcal{C}^1$, mais on ne parvient pas avec ce même théorème à montrer qu'elle est $\mathcal{C}^2$. On effectue alors une intégration par parties dans l’expression de $F'(x)$ (ou un changement de variable) et on parvient ensuite à montrer que $F'$ est $\mathcal{C}^1$, donc que $F$ est $\mathcal{C}^2$.
\end{proof}

\begin{ex}[Lafitte]{}
	Soit $a$ et $b$ deux réels strictement supérieurs à $1$. Calculer : $$\int_{0}^{\pi}\ln\left(\frac{b-\cos(x)}{a-\cos(x)}\right)dx$$
\end{ex}
\begin{proof}
	Écrire cette intégrale comme $I(b)-I(a)$ puis dériver la fonction $I$, calculer la valeur de cette dérivée à l'aide des règles de Bioche, puis intégrer pour obtenir la fonction $I$ à constante près. La détermination de cette constante est inutile puisqu'ensuite on s'intéresse à $I(b)-I(a)$.
\end{proof}

\section{Intégrales dépendant d'un paramètre complexe}
\begin{ex}[Moulin.IP.31]{$\Delta$ Fonction $\Gamma$ de la variable complexe}
	On note $\mathcal{P}=\{z\in\C\ : \ \text{Re}(z)>0\}$. On définit la fonction Gamma sur $\mathcal{P}$ par : $$\Gamma(z)=\int_{0}^{+\infty}t^{z-1}e^{-t}\ dt$$ Elle vérifie : 
	\begin{itemize}
		\item $\forall n\in\N,\ \Gamma(n+1)=n!$ ;
		\item $\Gamma$ est log-convexe ;
		\item $\Gamma$ est continue sur $\mathcal{P}$
	\end{itemize}
\end{ex}
\begin{proof}
	Pour le premier point, effectuer des intégrations par parties successives. Pour le troisième, dominer.
\end{proof}

\begin{ex}[Moulin.IP.32]{Prolongement de la fonction $\zeta$}
	\begin{enumerate}
		\item Montrer que la fonction $z\mapsto\displaystyle\int_{1}^{+\infty}\frac{\lfloor t\rfloor+1-t}{t^{z+1}}\ dt$ est définie et continue sur $\mathcal{P}$ (voir exercice précédent pour la définition).
		\item En déduire que $$\zeta:z\mapsto\sum_{n=1}^{+\infty}\frac{1}{n^z}$$ admet un prolongement continu à $\mathcal{P}\setminus\{1\}$.
	\end{enumerate}
\end{ex}

\begin{ex}[Guindy]{Une démonstration du théorème de D'Alembert-Gauss}
	\begin{enumerate}
		\item Soit $f$ une fonction de $\R$ dans $\C\st$, $2\pi$-périodique de classe $\mathcal{C}^1$. Montrer : $$\frac{1}{2i\pi}\int_{0}^{2\pi}\frac{f'(t)}{f(t)}dt\in\Z$$
		\item En déduire le théorème de D'Alembert-Gauss.
		\item Calculer la valeur de cette intégrale pour $P\in\C[X]$.
	\end{enumerate}
\end{ex}

\section{Théorèmes de Fubini}

\begin{ex}[Moulin.IP.33]{Théorème de Fubini sur les rectangles}
	Soit $[a,b]$ et $[c,d]$ deux segment non triviaux de $\R$ et $f:[a,b]\times[c,d]\to\C$ une fonction continue.
	\begin{enumerate}
		\item Montrer l'existence des intégrales $$\int_{a}^{b}\int_{c}^{d}f(x,y)\ dy\ dx \ \ \text{et} \ \ \int_{c}^{d}\int_{a}^{c}f(x,y)\ dx\ dy$$
		\item En considérant la fonction $$t\mapsto \int_{a}^{b}\left(\int_{c}^{t}f(x,y)\ dy\right)\ dx$$ montrer l'égalité $$\int_{a}^{b}\int_{c}^{d}f(x,y)\ dy\ dx =\int_{c}^{d}\int_{a}^{c}f(x,y)\ dx\ dy$$
	\end{enumerate}
\end{ex}
\begin{proof}
	\begin{enumerate}
		\item Domination facile car continue sur un compact.
		\item Dériver puis intégrer entre $c$ et $d$.
	\end{enumerate}
\end{proof}

\begin{ex}[Moulin.IP.34]{Théorème de Fubini sur les produits d'intervalles}
	Soit $I$ et $J$ deux intervalles non triviaux de $\R$ et $f:I\times J\to\R_+$ une fonction continue à valeurs positives. On fait les hypothèses suivantes : 
	\begin{enumerate}[label=(\roman*)]
		\item les fonctions $x\mapsto\displaystyle\int_J f(x,y)\ dy$ et $y\mapsto\displaystyle\int_I f(x,y)\ dy$ sont définies et continues par morceaux respectivement sur $I$ et $J$ ;
		\item l'intégrale $\displaystyle\int_I\left(\int_J f(x,y)\ dy\right)\ dx$ converge.
	\end{enumerate}
	\begin{enumerate}
		\item En utilisant le résultat de l'exercice précédent, montrer, pour tout segment $K$ inclus dans $J$, l'égalité : $$\int_K\left(\int_I f(x,y)\ dx\right)\ dy=\int_I\left(\int_K f(x,y)\ dy\right)\ dx$$
		\item En déduire la convergence de $\displaystyle\int_J\left(\int_I f(x,y)\ dx\right)\ dy$ et l'égalité : $$\int_I\left(\int_J f(x,y)\ dy\right)\ dx=\int_J\left(\int_I f(x,y)\ dx\right)\ dy$$
	\end{enumerate}
\end{ex}

\section{Méthode de Laplace}

\begin{ex}[Moulin.IP.35]{$\Delta$ Méthode de Laplace}
	Soit $f$ une fonction continue et à valeurs strictement positives sur $[0,a]$ avec $a>0$. on suppose qu'il existe $(\alpha,p)\in(\R_+\st)^2$ tel que $f(t)=1-\alpha t^p+o(t^p)$ au voisinage de $0$ et $\forall t\in]0,a],\ f(t)<1$. On se propose de trouver un équivalent de $$I(x)=\int_{0}^{a}f(t)^x\ dt$$ lorsque $x$ tend vers $+\infty$.
	\begin{enumerate}
		\item Montrer qu'il existe $\beta>0$ tel que $$\forall t\in[0,a],\ f(t)\leqslant\exp(-\beta t^p)$$
		\item Déterminer la limite, quand $x$ tend vers $+\infty$, de $$g(x)=\int_{0}^{ax^{1/p}}f\left(ux^{-1/p}\right)^x\ dt$$
		\item Conclure.
	\end{enumerate}
\end{ex}
\begin{proof}
	Bel exercice qui permet de déterminer rapidement des équivalents d'intégrales moches.
	\begin{enumerate}
		\item Au voisinage de $0$, utiliser des équivalents puis sur le reste (compact), utiliser le théorème des bornes atteintes pour bien choisir la constante. Prendre le minimum des deux constantes obtenues.
		\item On pose $$f(u,x)=\exp\left(x\ln f\left(ux^{-1/p}\right)\right)$$ si $u\in[0,ax^{-1/p}]$ et $0$ sinon. Alors $$f(u,x)\sim \exp(-\alpha u^p)=:h(u)$$ La question 1 donne la domination $$\lvert f(u,x)\rvert\leqslant\exp(-\beta u^{1/p})$$ Par passage à la limite sous l'intégrale, on obtient : $$\int_{0}^{ax^{1/p}}f\left(ux^{-1/p}\right)^x\ dt\xrightarrow[x\to+\infty]{}\int_{0}^{+\infty}e^{-\alpha u^p}\ du$$
		\item On fait le changement de variable $t=ux^{-1/p}$ dans la définition de $g$. Alors, la question $2$ fournit : $$\int_{0}^{a}f(t)^x\ dt\sim \frac{1}{x^{1/p}}\int_{0}^{+\infty}e^{-\alpha u^p}\ du$$
	\end{enumerate}
\end{proof}

\begin{ex}[Moulin.IP.36]{Application 1}
	Donner un équivalent, lorsque $\lambda$ tend vers $+\infty$ de $$I(\lambda)=\int_{-\pi/4}^{\pi/4}\exp(-\lambda\sin^2(x))\ dx$$
\end{ex}
\begin{proof}
	Utiliser la parité de l'intégrande pour se ramener à $[0,\pi/4]$. Ensuite, écrire : $$\exp(-\sin^2(x))=1-x^2+o(x^2)$$ et vérifier que cette fonction vérifie les hypothèses de la méthode de Laplace. On en déduit : $$I(\lambda)=\int_{-\pi/4}^{\pi/4}\exp(-\lambda\sin^2(x))\ dx\sim\frac{2}{\sqrt{\lambda}}\int_{0}^{+\infty}e^{-u^2}\ du=\sqrt{\frac{\pi}{\lambda}}$$
\end{proof}

\begin{ex}[Moulin.IP.37]{$\star$ Application 2}
	Convergence, limite et équivalent lorsque $n$ tend vers l'infini de la suite $(I_n)$ définie par $$I_n=\int_{0}^{+\infty}\frac{dt}{(1+t^3)^n}$$
\end{ex}
\begin{proof}
	La convergence de chaque $I_n$ se fait bien. la limite est nulle avec le théorème de convergence dominée (domination par le rang $n=1$).  Ensuite, découper $\R_+$ en $[0,2]$ et $[2,+\infty[$ et appliquer la méthode de Laplace sur le premier segment. Pour le deuxième intervalle, on montre aisément que l'intégrale est négligeable devant celle sur $[0,2]$. Or : $$\frac{1}{1+t^3}=1-t^3+o(t^3)$$ et vérifie les hypothèses. On en déduit finalement que $$I_n\sim\frac{1}{\sqrt[3]{n}}\int_{0}^{+\infty}e^{-u^3}\ du$$ On pourrait aussi obtenir cet équivalent avec le CDV $u=\sqrt[3]{n}t$ et un TCD, mais justifier la domination serait beaucoup plus embêtant car il faudrait étudier le signe de la dérivée de $$t\mapsto\left(1+\frac{u^3}{t}\right)^{-t}$$ à $u$ fixé et utiliser l'inégalité : $$\ln(1+t)\geqslant1-\frac{1}{1+t}$$
\end{proof}

\begin{ex}[Moulin.IP.38]{$\star$ Application 3}
	Déterminer un équivalent de $$I_n=\int_{0}^{1}(2-\cosh(t))^n\ dt \ \ \text{et} \ \ J_n=\int_{0}^{1}(2-\sinh(t))^n\ dt$$
\end{ex}
\begin{proof}
	Celui-ci est plus difficile.
	\begin{itemize}
		\item On a $$2-\cosh(t)=1-\frac{t^2}{2}+o(t^2)$$ et la fonction vérifie les hypothèses de la méthode de Laplace. On obtient avec cette méthode et après changement de variable $t=u/\sqrt{2}$ avec l'intégrale de Gauss l'équivalent $$I_n\sim\sqrt{\frac{\pi}{2n}}$$ Cela ressemble à du Wallis !
		\item Pour la deuxième, il faut diviser par $2$ pour se ramener aux hypothèses. On a : $$1-\frac{\sinh(t)}{2}=1-\frac{t}{2}+o(t)$$ et cela vérifie les hypothèses donc : $$\int_{0}^{1}\left(1-\frac{\sinh(t)}{2}\right)^n\ dt\sim\frac{1}{n}\int_{0}^{+\infty}e^{-t/2}\ dt=\frac{2}{n}$$ D'où le résultat final : $$J_n\sim\frac{2^{n+1}}{n}$$
	\end{itemize}
\end{proof}

\section{Transformation de Laplace}

\begin{ex}[Moulin.IP.12]{Transformée de Laplace}
	Soit $f\in\mathcal{C}([0,+\infty[,\C)$ ayant une limite $l\in\C$ en $+\infty$.
	\begin{enumerate}
		\item Montrer que pour tout $x>0$, la fonction $t\mapsto e^{-xt}f(t)$ est intégrable sur $\R_+$.
		\item Montrer que $\displaystyle\lim_{x\to 0} x\int_{0}^{+\infty}e^{-xt}f(t)\ dt=l$
		\item Montrer que $\displaystyle\lim_{x\to +\infty} x\int_{0}^{+\infty}e^{-xt}f(t)\ dt=f(0)$
		\item Existe-t-il $f\in\mathcal{C}([0,+\infty[,\C)$ telle que $x\mapsto\displaystyle x\int_{0}^{+\infty}e^{-xt}f(t)\ dt$ admette une limite finie en $0$ sans que $f$ n'ait de limite finie en $+\infty$ ?
	\end{enumerate}
\end{ex}
\begin{proof}
	Classique, du cours pour les SIistes.
	\begin{enumerate}
		\item Continue admet une limite donc bornée et on a le grand O qui suffit. Ensuite, on remarque par changement de variable que  : $$x\int_{0}^{+\infty}e^{-xt}f(t)\ dt=\int_{0}^{+\infty}e^{u}f\left(\frac{u}{x}\right)\ du$$ Et on a les continuités (par morceaux) réglementaires et la domination par une constante multipliée par $e^{-u}$.
		\item Passage à la limite sous le signe intégrale avec ce que l'on a fait avant.
		\item Passage à la limite sous le signe intégrale avec ce que l'on a fait avant.
		\item Considérer $t\mapsto e^{it}$ (penser aux transformées de Laplace de $\sin$ et $\cos$ qui existent en $SI$). 
	\end{enumerate}
\end{proof}

\begin{ex}[Châteaux]{Premières propriétés}
	Soit $f:\R_+\to\R$ continue par morceaux. Lorsque l'intégrale existe, on pose $$L(f)(p)=\int_{0}^{+\infty}f(t)e^{-pt}\ dt$$ la transformée de Laplace de $f$. On note $\mathcal{A}_f$ l'ensemble des $p\in\C$ tels que $t\mapsto e^{-pt}f(t)$ soit intégrable sur $]0,+\infty[$.
	\begin{enumerate}
		\item \begin{enumerate}[label=\alph*.]
			\item Si $f$ est intégrable sur $]0,+\infty[$, montrer que $L(f)$ existe, est définie et continue sur $\R_+\st$, est bornée et tend vers $0$ en $+\infty$.
			\item Montrer que si $f(t)=O(e^{p_0t})$ en $+\infty$, alors $L(f)$ est définie et continue sur $]p_0,+\infty[$.
		\end{enumerate}
		\item On suppose que $f$ est à croissance polynomiale, c'est-à-dire qu'il existe $k\in\N$ tel qu'en $+\infty$, $f(t)=O(t^k)$. Démontrer que $L(f)$ est de classe $\mathcal{C}^\infty$ sur $]0,+\infty[$ et $$(L(f))^{(n)}(p)=(-1)^n\int_{0}^{+\infty}t^nf(t)e^{-pt}\ dt$$ Démontrer que $L(f)$ et toutes ses dérivées tendent vers $0$ en $+\infty$.
		\item Soit $p_0\in\mathcal{A}_f$. Démontrer que tout nombre complexe $p$ tel que $\text{Re}(p)\geqslant\text{Re}(p_0)$ appartient à $\mathcal{A}_f$. On pose alors $\sigma_f=\inf \R\cap\mathcal{A}_f$.
		\item Soit $p_0$ tel que $L(f)(p_0)$ existe. Démontrer que, pour tout nombre complexe $p$ tel que $\text{Re}(p)>\text{Re}(p_0)$, $L(f)(p)$ existe.
	\end{enumerate}
\end{ex}

\begin{ex}[Châteaux]{Exemples de transformées de Laplace}
	Calculer $\sigma_f$ et $L(f)$ pour les fonctions suivantes : 
	\begin{enumerate}
		\item $f:t\mapsto t$ ;
		\item $f:t\mapsto\exp(\alpha t)$ ;
		\item $f:t\mapsto\cos(\omega t)$ ;
		\item $f:t\mapsto\sin(\omega t)$ ;
		\item $f:t\mapsto t^ne^{at}$.
	\end{enumerate}
\end{ex}

\begin{ex}[Châteaux]{Application au calcul d'une intégrale}
	On se propose de démontrer que $\displaystyle\int_{0}^{+\infty}\frac{\sin(t)}{t}\ dt=\frac{\pi}{2}$. On pose $G(p)=\displaystyle\int_{0}^{+\infty}e^{-pt}\frac{\sin(t)}{t}\ dt$.
	\begin{enumerate}
		\item Montrer que $G$ est de classe $\mathcal{C}^1$ sur $]0,+\infty[$ et que $G'(p)=\displaystyle\frac{-1}{p^2+1}$.
		\item En déduire que $G(p)=\displaystyle\frac{\pi}{2}-\arctan(p)$ pour $p>0$.
		\item Montrer que $G$ est continue en $0$. Indication : on posera $F(t)=\displaystyle-\int_{t}^{+\infty}\frac{\sin(u)}{u}\ du$ et on remarquera que $G(p)-G(0)=p\int_{0}^{+\infty}e^{-pt}F(t)\ dt$.
		\item Conclure.
	\end{enumerate}	
\end{ex}

\begin{ex}[Châteaux]{Théorèmes de la valeur initiale et finale}
	\begin{enumerate}
		\item Théorème de la valeur initiale : soit $f$ continue sur $\R_+\st$ de limite $l$ en $0$ et vérifiant $f(t)=O(e^{at})$ en $+\infty$. Montrer que $$pL(f)(p)\xrightarrow[p\to+\infty]{}l$$
		\item Théorème de la valeur finale : soit $f$ continue sur $\R_+$ de limite $l$ en $+\infty$. Montrer que $$pL(f)(p)\xrightarrow[p\to0^+]{}l$$
	\end{enumerate}
\end{ex}
\begin{ex}[Châteaux]{Applications}
	\begin{enumerate}
		\item Retrouver la valeur de $\displaystyle\int_{0}^{+\infty}\frac{\sin(t)}{t}\ dt$ en appliquant le théorème de la valeur finale avec $f(t)=-\displaystyle\int_{t}^{+\infty}\frac{\sin(u)}{u}\ du$.
		\item Déterminer un équivalent quand $p$ tend vers $+\infty$ de $\displaystyle\int_{1}^{+\infty}\frac{e^{-pt}}{t}\ dt$. En déduire un équivalent en $+\infty$ de la transformée de Laplace de $\displaystyle t\mapsto\frac{1}{t+1}$.
	\end{enumerate}
\end{ex}

\begin{ex}[Châteaux]{Transformée de Laplace des dérivées successives}
	Soit $f$ de classe $\mathcal{C}^1$
	telle que $L(f)$ et $L(f')$ existent sur $]p_0,+\infty[$, $f(0)=0$ et $f(t)=o(e^{pt})$ en $+\infty$ pour $p>p_0$. Montrer que $L(f')(p)=pL(f)(p)$. Généraliser.
\end{ex}

\begin{ex}[Châteaux]{Injectivité de la transformation de Laplace}
	Soit $f$ continue sur $\R_+$ telle que $f(t)=O(e^{p_0t})$ en $+\infty$. On suppose que $L(f)=0$ sur $]p_0,+\infty[$. Montrer que $f=0$.
\end{ex}

\section{Convolution}

\begin{ex}[Châteaux]{Propriétés de base de la convolution sur $\R$}
	Si $f$ et $g$ sont des fonctions continues par morceaux de $\R$ dans $\R$, on appelle produit de convolution de $f$ et $g$, $f\star g$, la fonction, si elle existe, définie par $$f\star g:x\mapsto\displaystyle\int_\R f(x-t)g(t)\ dt$$Le support d'une fonction $f$ est l'ensemble définie par $$\text{Supp}(f)=\text{Adh}\left\{x\in\R\ \mid \ f(x)\ne 0\right\}$$
	\begin{enumerate}
		\item Montrer que si $f$ et $g$ sont à support compact, alors $f\star g$ est aussi à support compact.
		\item Démontrer que le produit de convolution est commutatif. On admet qu'il est associatif.
		\item Montrer que si $f$ est bornée et $g$ est intégrable, alors $f\star g$ existe et est bornée par $\lVert f\rVert_\infty\lVert g\rVert_1$.
		\item Montrer que si $f$ et $g$ sont de carré intégrable, alors $f\star g$ existe et est bornée par $\lVert f\rVert_2\lVert g\rVert_2$.
		\item Montrer que si $f$ est continue et bornée et si $g$ est intégrable, alors $f\star g$ est continue. Remarque : la continuité reste vraie en supposant uniquement les fonctions continues par morceaux.
		\item Montrer que si $f$ est de classe $\mathcal{C}^1$ à dérivée bornée et si $g$ est intégrable, alors $f\star g$ est de classe $\mathcal{C}^1$ et $(f\star g)'=f'\star g$.
	\end{enumerate}
\end{ex}

\begin{ex}[Châteaux]{Exemples de convolution}
	\begin{enumerate}
		\item Pour $a>0$, on note $\varphi_a$ la fonction $t\mapsto\exp(-a^2t^2)$. Montrer que $$\varphi_a\star\varphi_b=\sqrt{\frac{\pi}{a^2+b^2}}\varphi_{\frac{ab}{\sqrt{a^2+b^2}}}$$
		\item On note $\Pi_T$ la fonction nulle à l'extérieur de $\displaystyle\left[-\frac{T}{2},\frac{T}{2}\right]$ et égale à $1$ sur $\displaystyle\left[-\frac{T}{2},\frac{T}{2}\right]$. Calculer $\Pi_T\star\Pi_T$.
		\item On note $e_a:t\mapsto e^{-a\lvert t\rvert}$. Calculer $e_a\star e_b$.
	\end{enumerate}
\end{ex}

\begin{ex}[Châteaux]{Unités approchées}
	\begin{enumerate}
		\item Soit $f$ continue bornée sur $\R$. Soit $(\varphi_n)$ une suite de fonctions continues positives intégrables telle que $\displaystyle\int_\R\varphi_n(t)\ dt\xrightarrow[n\to+\infty]{}1$ et, pour tout $\delta>0$, $\displaystyle\int_{\lvert t\rvert\geqslant \delta}\varphi_n(t)\ dt\xrightarrow[n\to+\infty]{}0$. Montrer que la suite de fonctions $(f\star\varphi_n)$ converge simplement vers $f$, et uniformément sur tout segment.
		\item Exemples : on peut par exemple prendre $$\varphi_n(t)=\frac{1}{\pi}\frac{n}{1+n^2t^2}\ \ \text{ou}\ \ \varphi_n(t)=\sqrt{\frac{n}{\pi}}e^{-nt^2}$$
	\end{enumerate}
\end{ex}

\begin{ex}[Châteaux]{Application de la convolution à l'approximation}
	Pour $n\in\N$, on pose $\displaystyle g_n(x)=\frac{1}{c_n}(1-x^2)^n$ si $x\in[-1,1]$ avec $c_n=\displaystyle\int_{-1}^{1}(1-t^2)^n\ dt$. Soit $f:[-1,1]\to\R$ nulle en dehors de $\displaystyle\left[-\frac{1}{2},\frac{1}{2}\right]$.
	\begin{enumerate}
		\item Montrer que la restriction de $f\star g_n$ à $\displaystyle\left[-\frac{1}{4},\frac{1}{4}\right]$ est polynomiale.
		\item Montrer que la suite de fonctions $(f\star g_n)$ converge uniformément vers $f$ sur $\displaystyle\left[-\frac{1}{4},\frac{1}{4}\right]$.
		\item En déduire une preuve du théorème de Weierstrass.
	\end{enumerate}
\end{ex}

\begin{ex}[Châteaux]{Espace de Schwartz et transformation de Fourier}
	On note $\mathcal{S}$ l'espace de Schwartz, c'est-à-dire de classe $\mathcal{C}^\infty$ et dont toutes les dérivées sont des $O(t^k)$ à l'infini pour tout $k\in\N$. Si $f\in\mathcal{S}$, on note $\widehat{f}$ sa transformée de Fourier, c'est-à-dire la fonction : $$\widehat{f}:\omega\mapsto\int_\R f(t)e^{-i\omega t}\ dt$$ Montrer que si $f$ et $g$ appartiennent à $\mathcal{S}$, alors $f\star g$ appartient à $\mathcal{S}$ et $$\widehat{f\star g}=\widehat{f}\times\widehat{g}$$ On pourra s'autoriser à échanger des intégrales doubles sans justification.
\end{ex}



\end{document}